% Created 2026-03-17 ter 14:20
% Intended LaTeX compiler: lualatex
\documentclass[12pt]{article}
\usepackage[a4paper,left=3cm,top=3cm,right=2cm,bottom=2cm]{geometry}
\setcounter{secnumdepth}{4}
\setcounter{tocdepth}{2}
\usepackage{graphicx}
\usepackage{longtable}
\usepackage{wrapfig}
\usepackage{rotating}
\usepackage[normalem]{ulem}
\usepackage{amsmath}
\usepackage{amssymb}
\usepackage{capt-of}
\usepackage{hyperref}
\usepackage{fgv-header}
\author{Gustavo M. Mendes de Tarso}
\date{\today}
\title{}
\hypersetup{
 pdfauthor={Gustavo M. Mendes de Tarso},
 pdftitle={},
 pdfkeywords={},
 pdfsubject={},
 pdfcreator={Emacs 28.2 (Org mode 9.5.5)}, 
 pdflang={English}}
\begin{document}

\begingroup
\linespread{1}\selectfont
\begin{tcolorbox}[title=Ficha Técnica,
  colback=gray!5,colframe=gray!40,boxrule=0.4pt,sharp corners]
\begin{tblr}{rowsep=1pt,stretch=1, rows={t}, colspec={Q[l,2.8cm] X[l]}}
\textbf{Curso}                   & Mestrado de Políticas Públicas e Governo \\
\textbf{Turma}                   & T-01 \\
\textbf{Pólo}                    & Brasília \\
\textbf{Disciplina}              & Teorias de Administração Pública \\
\textbf{Professor}               & Bernardo Buta \\
\textbf{Trabalho}                & Aplicação do protocolo PRISMA 2020 \\
\textbf{Aluno(s)}                & Gustavo M. Mendes de Tarso \\
\textbf{Título do trabalho}      & Seleção e síntese de artigo sobre governança pública, compliance e auditoria interna \\
\end{tblr}
\end{tcolorbox}
\endgroup

\FGVNewPage

\section{Introdução}
\label{sec:org57d27c0}

Este trabalho apresenta, em formato alinhado ao \textbf{PRISMA 2020}, o fluxo de identificação, triagem, elegibilidade e inclusão de estudos relacionados aos temas governança pública, compliance e administração pública. A estratégia de busca foi aplicada inicialmente na \textbf{Scopus} e, em seguida, na \textbf{Web of Science}, utilizando combinação de termos voltados a governança pública, compliance, integridade e revisão sistemática.

Na \textbf{Scopus}, a busca retornou 2 registros. Na \textbf{Web of Science}, a mesma lógica de busca retornou 100 registros, posteriormente refinados por \textbf{ano de publicação a partir de 2023} e por \textbf{idioma português}. Após esse refinamento, apenas 2 estudos permaneceram aderentes para compor a triagem final em conjunto com os resultados da Scopus.

\section{Estratégia de busca}
\label{sec:org1118687}

\subsection{Query utilizada}
\label{sec:org684f306}

\begin{quote}
TITLE-ABS-KEY("public governance" OR "public sector governance" OR "governança pública")
AND TITLE-ABS-KEY(compliance OR integrity OR "public integrity")
AND TITLE-ABS-KEY("systematic review" OR "systematic literature review")
\end{quote}

\subsection{Bases consultadas}
\label{sec:org9307204}
\begin{itemize}
\item Scopus
\item Web of Science
\end{itemize}

\subsection{Refinamentos aplicados}
\label{sec:org3374d15}
\begin{itemize}
\item Web of Science: recorte temporal a partir de 2023
\item Web of Science: idioma português
\end{itemize}

\section{PRISMA 2020}
\label{sec:org33dd295}

\subsection{Fluxo textual dos estudos}
\label{sec:org7e6aa5f}

\subsubsection{Identificação}
\label{sec:orga025a68}
\begin{itemize}
\item Registros identificados em bases de dados: \textbf{n = 102}
\begin{itemize}
\item Scopus: \textbf{n = 2}
\item Web of Science: \textbf{n = 100}
\end{itemize}

\item Registros removidos antes da triagem: \textbf{n = 98}
\begin{itemize}
\item Registros excluídos por refinamento de ano e idioma na Web of Science: \textbf{n = 98}
\item Duplicatas removidas: \textbf{n = 0}
\item Registros removidos por outros motivos: \textbf{n = 0}
\end{itemize}
\end{itemize}

\subsubsection{Triagem}
\label{sec:orgc016fd8}
\begin{itemize}
\item Registros triados: \textbf{n = 4}
\item Registros excluídos: \textbf{n = 2}
\begin{itemize}
\item Tema não aderente ao objeto da pesquisa: \textbf{n = 1}
\item Estudo não caracterizado como revisão sistemática aderente ao recorte temático: \textbf{n = 1}
\end{itemize}
\end{itemize}

\subsubsection{Elegibilidade}
\label{sec:org832b977}
\begin{itemize}
\item Relatórios buscados para recuperação: \textbf{n = 2}
\item Relatórios não recuperados: \textbf{n = 1}
\begin{itemize}
\item Motivo: indisponibilidade de acesso ao texto completo do estudo sobre blockchain
\end{itemize}

\item Relatórios avaliados para elegibilidade: \textbf{n = 1}
\item Relatórios excluídos após leitura em texto completo: \textbf{n = 0}
\end{itemize}

\subsubsection{Inclusão}
\label{sec:orgaeb5cc3}
\begin{itemize}
\item Estudos incluídos na revisão: \textbf{n = 1}
\end{itemize}

\subsection{Tabela-resumo do fluxo PRISMA 2020}
\label{sec:org97838b3}

\begin{table}[htbp]
\centering
\scriptsize
\setlength{\tabcolsep}{4pt}
\renewcommand{\arraystretch}{1.15}
\begin{tabularx}{\textwidth}{>{\raggedright\arraybackslash}X >{\raggedleft\arraybackslash}m{0.18\textwidth}}
\toprule
Etapa PRISMA 2020 & n \\
\midrule
Registros identificados em bases de dados & 102 \\
Registros removidos antes da triagem & 98 \\
Registros triados & 4 \\
Registros excluídos & 2 \\
Relatórios buscados para recuperação & 2 \\
Relatórios não recuperados & 1 \\
Relatórios avaliados para elegibilidade & 1 \\
Relatórios excluídos após leitura completa & 0 \\
Estudos incluídos na revisão & 1 \\
\bottomrule
\end{tabularx}
\normalsize
\end{table}

\subsection{Estudos classificados no fluxo}
\label{sec:org4c86d1a}

\begin{table}[htbp]
\centering
\scriptsize
\setlength{\tabcolsep}{3pt}
\renewcommand{\arraystretch}{1.15}
\begin{tabularx}{\textwidth}{>{\raggedright\arraybackslash}X >{\raggedright\arraybackslash}m{0.17\textwidth} >{\raggedright\arraybackslash}m{0.22\textwidth} >{\raggedright\arraybackslash}m{0.25\textwidth}}
\toprule
Estudo & Base & Situação no fluxo & Motivo \\
\midrule
Exploring the internal audit of public procurement governance: a systematic literature review & Scopus & Incluído & Aderente ao tema e disponível \\
Blockchain Interoperability in Secure E-Governance: A Systematic Literature Review & Scopus & Não recuperado & Texto completo indisponível \\
Consensual Administration, Accountability and Transparency in the Brazilian Public Administration & Web of Science & Excluído na triagem & Não configurado como revisão sistemática aderente ao recorte específico \\
The Brazilian Paralympic Committee as a Civil Society Organization of Public Interest & Web of Science & Excluído na triagem & Tema não aderente ao objeto \\
\bottomrule
\end{tabularx}
\normalsize
\end{table}

\subsection{Representação esquemática em PlantUML}
\label{sec:org10b8ba6}

\begin{center}
\includesvg[width=.9\linewidth]{prisma2020_fluxo}
\end{center}

\section{Resumo do artigo selecionado}
\label{sec:org52283ac}

O artigo selecionado, intitulado \emph{Exploring the internal audit of public procurement governance: a systematic literature review}, examina o papel da auditoria interna na governança das compras públicas por meio de uma revisão sistemática da literatura. O estudo parte do entendimento de que as contratações públicas constituem uma área especialmente sensível a falhas de controle, desperdícios, riscos operacionais e práticas de corrupção, razão pela qual mecanismos de auditoria interna assumem papel central na promoção da transparência, da accountability e da eficiência administrativa.

Em termos metodológicos, o artigo adota a abordagem de \textbf{systematic literature review} com apoio explícito do referencial PRISMA. A revisão cobre publicações entre 2001 e 2024 e reúne 68 estudos, a partir dos quais os autores procuram mapear a evolução do campo, identificar os principais eixos temáticos, apontar limitações recorrentes e propor uma agenda futura de pesquisa. O trabalho demonstra que a produção científica sobre auditoria interna em compras públicas cresceu significativamente a partir de 2018, indicando maior preocupação internacional com integridade, rastreabilidade e conformidade nas aquisições governamentais.

Os resultados mostram que a auditoria interna é compreendida na literatura como instrumento decisivo para o fortalecimento da governança pública, especialmente por sua capacidade de prevenir fraudes, monitorar conformidade normativa, aperfeiçoar controles internos e reduzir vulnerabilidades nas contratações. O estudo também destaca que sua eficácia prática depende de fatores institucionais concretos, como independência funcional dos auditores, qualificação técnica, apoio organizacional, integração com gestão de riscos e acompanhamento efetivo das recomendações emitidas.

Outro ponto relevante do artigo é a identificação de novos vetores tecnológicos aplicados à auditoria e à governança de compras públicas, como análise de dados em larga escala, automação de processos e ferramentas digitais voltadas à detecção de irregularidades. Ainda assim, os autores assinalam que muitos órgãos públicos enfrentam obstáculos persistentes, tais como escassez de recursos, baixa capacidade institucional, deficiência de treinamento e limitações tecnológicas, o que reduz o potencial transformador da auditoria interna.

No plano teórico, o artigo sustenta que o campo ainda é relativamente dependente de abordagens tradicionais e propõe maior integração entre perspectivas como \textbf{institutional theory}, \textbf{stakeholder theory} e \textbf{stewardship theory}, de modo a compreender com mais profundidade a complexidade da governança pública. Em síntese, a conclusão do estudo é que a auditoria interna deve ser vista não apenas como instância de controle formal, mas como mecanismo estratégico de integridade, aprimoramento institucional e fortalecimento da governança nas compras públicas.

\section{Considerações finais}
\label{sec:org17b3a51}

À luz do fluxo PRISMA 2020 e do artigo efetivamente incluído, observa-se que o recorte de governança pública e compliance no setor público, quando associado ao critério de revisão sistemática, produz universo bibliográfico relativamente restrito após a aplicação dos filtros de elegibilidade. Ainda assim, o estudo selecionado oferece contribuição relevante para a disciplina \textbf{Teorias de Administração Pública}, sobretudo por evidenciar que governança, accountability, integridade e controle interno são dimensões interdependentes na administração pública contemporânea.
\end{document}